\section{Introduction}\label{sec:introduction}

A triple of non-Archimedean norms contains more information than the
three relative spectral lists. The harmonic-path problem of
Haiden--Katzarkov--Kontsevich--Pandit asks for a geometric reconstruction
that retains this additional information: a path of non-expansive norms
whose spectral velocities trace a labelled network in the plane
\cite[\S\S7.2--7.3]{HKKP26}. We construct such paths for the identity
representation of the $A_2$ quiver, for every full boundary triple and in
every finite rank. The principal point is finite algebraic regularity of
a variational minimum in a completed norm space. A separate converse
identifies that minimum among paths satisfying an explicit first-order
matching condition.

\subsection{The boundary problem}
Let $K$ be the real-exponent Novikov field over the trivially valued field
$\CC$. Thus its elements are formal sums
\[
 f=\sum_{\lambda\in\RR}a_\lambda T^\lambda
\]
whose support is bounded below and has finite intersection with
$(-\infty,b]$ for every $b\in\RR$. For $f\ne0$ set
$\nu(f)=\min\{\lambda:a_\lambda\ne0\}$ and $|f|=e^{-\nu(f)}$.
We write $\OO=\{|f|\le1\}$ and $\mm=\{|f|<1\}$; the residue field is
$\OO/\mm=\CC$. This field is complete but not spherically complete.
Its spherically complete Hahn extension will be used in the proof; the
path constructed by the theorem is defined over $K$.

Fix $V=K^n$ and consider the quiver with one arrow from its vertex $C$
to its vertex $A$. We use the representations
\begin{equation}\label{eq:ABC}
 A=(0\longrightarrow K),\qquad
 B=(K\xrightarrow{\id}K),\qquad
 C=(K\longrightarrow0),\qquad E=B^n,
\end{equation}
and the central charges
\begin{equation}\label{eq:charges}
 Z(A)=e^{\pi i/6},\qquad Z(B)=i,\qquad Z(C)=e^{5\pi i/6}.
\end{equation}
A norm $h$ on $V$ is \emph{diagonalizable}, or \emph{split}, if there are
a basis $e_1,\ldots,e_n$ and real weights $w_1,\ldots,w_n$ such that
\begin{equation}\label{eq:splitnorm}
 h\Bigl(\sum_i b_i e_i\Bigr)=\max_i e^{w_i}|b_i|.
\end{equation}
The norms diagonalized by a fixed basis form an apartment. We retain their
scales; norms are not taken modulo multiplication by a positive scalar.
Two split norms admit a common orthogonal basis. Their $d_2$ distance
is the Euclidean norm of the differences of their logarithmic weights
in such a basis. In particular, this is the unnormalized metric, including
the determinant direction.

A norm on $E$ is a pair $c=(c_A,c_C)$. The identity arrow is non-expansive
exactly when
\begin{equation}\label{eq:order}
 c_A\le c_C
\end{equation}
pointwise on $V$. For a split norm $h$, put
$\overline V_h=\{h\le1\}/\{h<1\}$. A non-expansive pair induces an
arrow $\overline V_{c_C}\to\overline V_{c_A}$ over $\CC$.
The harmonicity condition for a norm in \cite[Definition~7.3]{HKKP26}
is semistability of this reduced representation. Every representation
of this one-arrow quiver over $\CC$ is a direct sum of $A$, $B$ and $C$.
Their phases are distinct, so the nonzero semistable objects are precisely
$A^m$, $B^m$ and $C^m$.

The boundary datum in this paper is an arbitrary triple of split norms
\begin{equation}\label{eq:boundarytriple}
 (h_B,h_A^+,h_C^+)\quad\text{on the fixed space }V.
\end{equation}
It determines the trivial negative-end filtration of $E$ and the
positive-end filtration $0\subset A^n\subset E$, with quotient $C^n$.
Indeed, a non-expansive norm on $B^n$ is harmonic precisely when its
reduced arrow is invertible. In a common orthogonal basis for its two
norms, this says that every relative weight is zero, hence that the two
norms agree. Thus $h_B$ is exactly a harmonic norm on the negative-end
object, while $h_A^+$ and $h_C^+$ are independent harmonic norms on the
two positive-end factors. No initial comparison between these three
norms is required.

We fix the area form and its transverse coordinate
\begin{equation}\label{eq:area}
 \omega=\frac{\dd x\wedge\dd t}{1+x^2+t^2},\qquad
 \lambda(x,t)=\frac{1}{\sqrt{1+t^2}}
                  \arctan\frac{x}{\sqrt{1+t^2}}.
\end{equation}
Let
\begin{equation}\label{eq:utime}
 \begin{gathered}
 a=\sqrt3,\qquad t=\sinh u,\qquad R(u)=\cosh u,\\
 \beta(u)=\arctan(a\tanh u),\qquad F(u)=\int_0^u\beta(s)\,\dd s.
 \end{gathered}
\end{equation}
We require the following limits in $d_2$:
\begin{equation}\label{eq:normalizedends}
 \begin{aligned}
 c_A(u),c_C(u)&\longrightarrow h_B &&(u\longrightarrow-\infty),\\
 e^{F(u)}c_A(u)&\longrightarrow h_A^+ &&(u\longrightarrow+\infty),\\
 e^{-F(u)}c_C(u)&\longrightarrow h_C^+ &&(u\longrightarrow+\infty).
 \end{aligned}
\end{equation}
These are the HKKP normalizations. In fact, for
$A_\mu(t)=\int_0^t\lambda(\mu s,s)\,\dd s$, one has
$A_a(\sinh u)=F(u)$ and $A_{-a}=-A_a$. The conjugate-charge direction
convention in \cite[Definition~7.8]{HKKP26} gives slopes $-a,0,a$
for the labels $A,B,C$, respectively.

\subsection{First-order matching and the main theorem}
We specify the comparison class in the uniqueness statement. This is
necessary because edge directions and abstract edge labels do not by
themselves identify the embeddings of the velocity subspaces at a join.

\begin{definition}\label{def:regularharmonic}
A \emph{finite regular harmonic path} is a continuous ordered pair of
split norms on the whole $u$-line with the following properties.
There is a finite subdivision into intervals on each of which both norms
lie in fixed apartments, with weight functions smooth up to the finite
endpoints. After subdivision at the intersections of the support branches,
the weights are grouped by their derivative functions
$p_g(u)\in(-\pi/2,\pi/2)$, and
\begin{equation}\label{eq:branch}
 R(u)\tan p_g(u)=\alpha_g\sinh u+b_g,
 \qquad \alpha_g\in\{-a,0,a\},\quad b_g\in\RR.
\end{equation}
Within each group the relative weights and reduced arrow are constant.
The actual reduced object is semistable, of type $A^{m_g}$, $B^{m_g}$ or
$C^{m_g}$, according as $\alpha_g=-a,0,a$. The closure of these labelled
branches is the support of the path.

The path satisfies \emph{first-order matching} if, at every finite
subdivision time $u_0$ and for $v=A,C$,
\begin{equation}\label{eq:matching}
 d_2\bigl(c_v(u_0),\midpt(c_v(u_0-h),c_v(u_0+h))\bigr)=o(h)
 \qquad(h\downarrow0).
\end{equation}
Here $\midpt$ is the midpoint of the canonical norm geodesic. This
condition is independent of bases and of artificial subdivisions.
\end{definition}

In a single apartment, \eqref{eq:matching} is equality of the two
one-sided first derivatives. More generally, it requires the one-sided
norm motions to fit to first order in the metric space of norms.
The groups in Definition~\ref{def:regularharmonic} retain their actual
reduced arrows, not merely their dimensions. In particular, a group with
both source and target is of type $B^{m_g}$ only when its reduced arrow
is an isomorphism.

\begin{theorem}\label{thm:main}
Let $n\ge1$. For every triple \eqref{eq:boundarytriple} of split norms
over $K$, there is a unique finite regular harmonic path satisfying
first-order matching and the normalized ends \eqref{eq:normalizedends}.
It has the following properties.
\begin{enumerate}
\item Its support is a finite balanced weighted honeycomb with incoming
$B$-multiplicity $n$ and outgoing $A$- and $C$-multiplicities $n$ each.
It has at most $n^2$ essential geometric vertices.
\item At every support vertex there is an actual reduced quiver
representation with two increasing-phase filtrations. Their successive
quotients are the incident incoming and outgoing edge objects. These
identifications agree along the edges and remove artificial subdivisions.
\item Each tail lies in fixed $K$-apartments. The two norms agree for
all sufficiently negative times. All three normalized limits have
$d_2$ error $O(|t|^{-1})$.
\item The path is equivariant under simultaneous $\GL_n(K)$ change
of basis of the three boundary norms.
\end{enumerate}
The path has analytic one-sided weights and satisfies the stronger
$O(h^2)$ estimate in \eqref{eq:matching}.
\end{theorem}

The multiplicities in the theorem specify the degree of the honeycomb;
they do not prescribe the valence of its vertices. In particular, crossing
vertices can be four-valent. The $n^2$ count is the standard finite
honeycomb bound from the finite dual subdivision in Knutson--Tao theory
\cite[Lemmas~4 and 6--8 and Theorem~1]{KT99}. The analytic content
is that the variational construction has finite support to which this
bound applies, and then a finite apartment description over the original
Novikov field. The theorem includes repeated boundary spectra,
coincident edges and collapsed bounded edges. Its estimates need not
be uniform over all boundary triples.

The uniqueness is equality of the actual norms at every time. It is
stronger than equality of support or S-equivalence within the stated
class. The theorem does not claim that every honeycomb with the same
boundary spacings gives the same path. Nor does the inclusion of degenerate
inputs by itself assert continuous dependence on arbitrary boundary data.
An explicit rank-three family below has a bounded face that collapses,
with a uniform Lipschitz estimate for its full norm paths.

\subsection{Relation to the harmonic-path conjecture}
HKKP conjecture a correspondence between harmonic paths and pairs of
ascending-phase filtrations with harmonic norms on their quotients
\cite[Conjecture~7.10]{HKKP26}. Their honeycomb example identifies the
boundary data with triples of diagonalizable norms
\cite[\S7.3]{HKKP26}. Theorem~\ref{thm:main} gives existence for every
such triple in the present field and area form. Its uniqueness statement
uses the finite regular class and the explicit condition
\eqref{eq:matching}.

The displayed conditions of \cite[Definition~7.8]{HKKP26} do not explicitly
impose this matching. They permit coordinate velocity subspaces to switch
at finitely many times while the support and abstract edge labels remain
unchanged. We give two rank-three examples with the same full normalized
ends and the same connected balanced support that differ by such switches.
They fail \eqref{eq:matching}. This explains why we do not state
Theorem~\ref{thm:main} as an unqualified proof of Conjecture~7.10.
The construction also supplies the local residue-category vertex data,
but we do not prove that abstract vertex filtrations alone imply
\eqref{eq:matching}, or identify every globally framed categorical network
with the norm paths in the theorem.

Haiden--Katzarkov--Simpson construct the category of spectral networks with
vertex filtrations and edge identifications
\cite[\S4.4]{HKS24}. In their constant-coefficient, product-type setting, the uniqueness
theorem gives the same support and S-equivalent reductions for networks
representing the same object after Novikov-field extension
\cite[Theorem~4.15]{HKS24}. Their regular-hexagon theorem supplies spectral
network representatives in a specified $A_2$ setting
\cite[Theorem~6.1]{HKS24}. These are different assertions from reconstruction
of the full framed norm path from arbitrary boundary norms. The family-Floer
comparison discussed in \cite[Remark~7.9]{HKKP26} is not used in our
existence or uniqueness proof.

The relation between honeycombs and feasible spectral triples is classical
\cite{KT99}; Moitra--Postnikov--Woodruff give an axiomatic treatment of
that relation \cite{MPW26}. Appleby--Whitehead construct real-valued Littlewood--Richardson
fillings from matrix pairs over valuation rings and hives from pairs
of full lattices \cite{AW13,AW18}. These constructions also retain
information beyond pairwise spectra. Our claim concerns the harmonic norm
path for the specified area form, rather than a first construction of
numerical honeycombs from valued linear algebra. Ben-Bassat's split
attainment theorem \cite[Theorem~1.1]{BB26} concerns the stationary harmonic-norm
potential of HKKP. Split attainment at each individual time does not supply
a coherent path or finite apartment structure; those require the
regularity argument below.

\subsection{The variational argument}
HKKP suggest a variational approach based on a convex function of the
velocity \cite[p.~85]{HKKP26}. For \eqref{eq:area}, its scalar kinetic
term in the $u$ coordinate is
\[
 \ell(v)=-\log\cos v,\qquad |v|<\pi/2.
\]
In a product apartment the local action is
\begin{equation}\label{eq:introlocalaction}
 \int_I R(u)\left\{\frac12\sum_{v=A,C}\sum_i\ell(w_{vi}'(u))
       +\frac a2\bigl(d(c_C(u))-d(c_A(u))\bigr)\right\}\dd u,
\end{equation}
where $d(h)$ is the logarithmic determinant norm in a common fixed volume
normalization. The positive-end reference velocities are
$-\beta$ and $\beta$. Subtracting their supporting affine terms gives a
nonnegative renormalized action on the whole line. A weighted strict
convexity estimate makes every minimizing sequence uniformly Cauchy.
The limit is first constructed in the complete space of all norms.

The main difficulty is to prove that this limit is a finite algebraic
path. Short ordered replacements give quadratic approximation by endpoint
geodesics and locally Lipschitz spectral velocities. Where groups of
velocities are separated, equality in the relevant Ky Fan inequalities
forces fixed orthogonal subspaces. At almost every time, scalar
normalization on these subspaces gives second-order norm approximations.
An endpoint repair with action error $o(h^3)$ permits variations by actual
residue subrepresentations. The resulting equations force semistability
and the three honeycomb directions.

The direction laws imply that each opposite-type pair of spectral
positions has at most one contact interval. There are consequently only
finitely many support events. On every interval between them the fixed
subspaces have scalar norm motion, giving a single apartment. A finite
valued-field transfer preserves the full endpoint norms, the reduced
arrows and the action, and identifies the resulting $K$-path with the
original minimum. Finally, opposite momentum directions at a matched
join give nonnegative Busemann boundary terms. This proves the converse
across all joins and yields the equality uniqueness in
Theorem~\ref{thm:main}.
